% =====================================================================
% Wave-13 A2: Kazhdan-voice adversarial assault on Lorgat 2020 §3,
% the isomorphism Sp_4(Z)/{\pm I_5} \simeq O(\Lambda^{3,2})_+/{\pm I_5}
% via the exterior-square / Pfaffian construction.
%
% Five ATTACK <-> HEAL cycles. Primary sources: Lorgat 2020 §3;
% Knapp, Lie Groups Beyond an Introduction (2002) VII.11; Helgason,
% Differential Geometry, Lie Groups, and Symmetric Spaces (1978) IX.5;
% Bourbaki, Groupes et algebres de Lie IV-VI planche IV (D_3 = A_3);
% Fulton-Harris Representation Theory Thm 19.20 (half-spin at signature
% (3,2)); Chevalley, Theory of Spinors (1954) II.4-II.7; Borel
% Introduction aux groupes arithmetiques 1969 §7 (Z-forms via lattices);
% Gritsenko-Nikulin 1995 Part I (lattice Lambda^{3,2} and the Siegel
% parabolic); Borcherds 1998 Invent 132 (cusp structure of O_+(L)).
%
% Chriss-Ginzburg voice. Subscripted kappa discipline (kappa not load-
% bearing here -- §3 is a group-theoretic identity). No AI attribution.
% Authorship: Raeez Lorgat.
% =====================================================================

\providecommand{\Sp}{\mathrm{Sp}}
\providecommand{\SL}{\mathrm{SL}}
\providecommand{\GL}{\mathrm{GL}}
\providecommand{\OO}{\mathrm{O}}
\providecommand{\SO}{\mathrm{SO}}
\providecommand{\Spin}{\mathrm{Spin}}
\providecommand{\PSp}{\mathrm{PSp}}
\providecommand{\Hb}{\mathbb{H}}
\providecommand{\Z}{\mathbb{Z}}
\providecommand{\Q}{\mathbb{Q}}
\providecommand{\R}{\mathbb{R}}
\providecommand{\C}{\mathbb{C}}
\providecommand{\sl}[1]{\mathfrak{sl}_{#1}}
\providecommand{\so}[1]{\mathfrak{so}_{#1}}
\providecommand{\sph}[1]{\mathfrak{sp}_{#1}}
\providecommand{\fg}{\mathfrak{g}}
\providecommand{\Pf}{\mathrm{Pf}}
\providecommand{\Gr}{\mathrm{Gr}}
\providecommand{\ClaimStatusTheorem}{\textbf{[Thm]}}
\providecommand{\ClaimStatusDerivation}{\textbf{[derivable]}}
\providecommand{\ClaimStatusConjectured}{\textbf{[Conj]}}
\providecommand{\ClaimStatusDefinition}{\textbf{[Def]}}

\section*{Wave 13 A2: Kazhdan-voice assault on Lorgat 2020 \S3 --- the
exterior-square isomorphism $\Sp_4(\Z)/\{\pm I\}\simeq\OO(\Lambda^{3,2})_+/\{\pm I\}$}
\label{sec:wave13-a2-sp4-orthogonal-kazhdan}

\noindent\textbf{Status.} \ClaimStatusTheorem.\ The isomorphism is a
theorem; it is the integral shadow of the accidental low-rank identity
$D_3\simeq A_3$ at the Lie-algebra level, $\sl{4}\simeq\so{3,2}$ via the
exterior square acting on $\Lambda^2\Q^4$. The Pfaffian construction
is explicit; the kernel is explicit; the half-spin sign picks out the
connected component indexed by $+$.

The mission: \emph{is} it a theorem? Five cycles stress-test every
move --- lattice integrality, sign conventions for $+$, accidentality at
$\Z$-form, parabolic compatibility, Weyl-group alignment. Every claim
is false until traced to a first-principles argument.

%---------------------------------------------------------------------
\subsection*{Cycle 1: the accidental isomorphism at the Lie-algebra level ---
$\sl{4}\simeq\so{3,2}$ via the exterior square}
\label{sec:w13a2-cycle-1}
%---------------------------------------------------------------------

\paragraph{Attack (Kazhdan voice).}
Lorgat 2020 \S3 (lines 233--350 of the paper text) states Lemma~1:
\[
\wedge^2:\Sp_4(\Z)/\{\pm I_5\}\;\xrightarrow{\ \sim\ }\;\SO^{+}(\Lambda^{3,2})\simeq\OO(\Lambda^{3,2})_+/\{\pm I_5\}.
\]
Before accepting the equals sign, the author owes four objections at
the arithmetic heart of the construction.

\begin{enumerate}[label=\textbf{O.\arabic*},leftmargin=1.5em]
\item \emph{Is this an isomorphism, or only a $\Z$-shadow of an
  $\R$-isogeny?} The Lie-algebra identity $\sl{4}\simeq\so{3,2}$
  (Dynkin type $A_3=D_3$) is well-known at semisimple complex Lie
  algebras (Bourbaki LIE IV--VI planche IV). At the real form
  $\sl{4,\R}\supset\Sp_4(\R)\cap\mathrm{real}$ versus
  $\so{3,2}\supset\OO(3,2)$, the isogeny is
  $\Spin(3,2)\twoheadrightarrow\SO^{+}(3,2)$ with kernel
  $\{\pm 1\}$; \emph{and} $\Sp_4(\R)\simeq\Spin(3,2)$ as real Lie
  groups (Helgason 1978 Ch.~IX.5; Knapp 2002 VII.11).
  The central isogeny $\Spin\to\SO$ has non-trivial fundamental
  group at each real form, and the passage from $\R$ to $\Z$ carries
  an additional index-$2$ or index-$4$ subtlety from the spin-norm.
  A blanket claim $\Sp_4(\Z)/\{\pm I\}\simeq\OO(\Lambda^{3,2})_+/\{\pm I\}$
  without an explicit $\Z$-level construction is mere $\R$-shadow.
\item \emph{Signature $(3,2)$ vs $(3,3)$ --- which one does $\wedge^2$
  land in?} Lorgat 2020 line 240 asserts that the pfaffian pairing
  $(u,v):=\text{coeff of }e_1\wedge e_2\wedge e_3\wedge e_4$ in
  $u\wedge v$, for $u,v\in\Lambda^2\Z^4$, is an even unimodular
  symmetric bilinear form of signature $(3,3)$. But the claimed
  orthogonal complement $\Lambda^{3,2}=(e_1\wedge e_3+e_2\wedge e_4)^\perp$
  has signature $(3,2)$. A signature-drop by $(0,1)$ when passing to
  an orthogonal hyperplane is exactly what one expects when the
  polarisation $q=e_1\wedge e_3+e_2\wedge e_4$ has negative square
  $(q,q)=-2$; that needs to be exhibited, not stipulated.
\item \emph{Which Pfaffian, which sign?} The form
  $(u,v):u\wedge v=(u,v)e_1\wedge e_2\wedge e_3\wedge e_4$ on
  $\Lambda^2\Z^4$ is antisymmetric under swap of the two
  $\Lambda^2\Z^4$-factors as $u\wedge v=-v\wedge u$; but its scalar
  value $(u,v)$ is \emph{symmetric}, because the swap is absorbed by
  the $\Z$-valued coefficient (swap $u\wedge v$, multiply by
  $(-1)^{4\cdot 4}=+1$). Verify this sign bookkeeping; many Pfaffian
  conventions differ by a factor of $2$ (Knuth, combinatorial
  identity involving $2^{n/2}$).
\item \emph{Integrality of the pfaffian pairing.} Is
  $(u,v)\in\Z$ for all $u,v\in\Lambda^2\Z^4$? The standard formula
  involves a denominator $n/2$ for $n\times n$ antisymmetric
  matrices; at $n=4$ this is $1/2$. The author owes an integrality
  argument, not a tacit assumption.
\end{enumerate}

\paragraph{Ghost theorem (first-principles reconstruction).}
The Dynkin accident $A_3=D_3$ means the root systems of $\sl{4}$ and
$\so{6}=\so{3,3}_\C$ coincide. Explicitly: let $V=\C^4$ with basis
$e_1,\ldots,e_4$; let $W=\Lambda^2 V$, a $6$-dimensional vector space.
The \emph{exterior square} representation
\[
\wedge^2:\GL(V)\longrightarrow\GL(W),\qquad g\mapsto (v_1\wedge v_2\mapsto gv_1\wedge gv_2),
\]
carries $\SL(V)$ into $\GL(W)$ preserving the natural
\emph{Pfaffian pairing}
\[
(\,,\,)_{\Pf}: W\otimes W\to \Lambda^4 V\simeq\C,\qquad u\otimes v\mapsto u\wedge v,
\]
where the identification $\Lambda^4 V\simeq\C$ is fixed by
$e_1\wedge e_2\wedge e_3\wedge e_4\leftrightarrow 1$. The pairing is
\emph{symmetric}: $u\wedge v=(-1)^{\deg u\cdot\deg v}v\wedge u=v\wedge u$
for $u,v\in W=\Lambda^2 V$ of even degree. It is
\emph{non-degenerate} with determinant $\pm 1$ in the basis
$\{e_i\wedge e_j\}_{i<j}$.

\emph{Computation of the signature.} In the basis
$e_1\wedge e_2,\ e_3\wedge e_4,\ e_1\wedge e_3,\ e_2\wedge e_4,\ e_1\wedge e_4,\ e_2\wedge e_3$
of $W$, the Gram matrix of $(\,,\,)_{\Pf}$ is
\[
G=\begin{pmatrix}0&1&&&&\\1&0&&&&\\&&0&1&&\\&&1&0&&\\&&&&0&-1\\&&&&-1&0\end{pmatrix},
\]
i.e.\ three hyperbolic planes with the third twisted by $-1$. Over
$\Z$ this is $\Lambda^{1,1}\oplus\Lambda^{1,1}\oplus\Lambda^{1,1}$;
over $\R$ the signature is $(3,3)$ (each hyperbolic plane contributes
$(1,1)$). This is precisely Lorgat 2020 line 241's claim, verified.

\emph{Symplectic reduction to signature $(3,2)$.} The symplectic
form $\omega=e_1^{\vee}\wedge e_3^{\vee}+e_2^{\vee}\wedge e_4^{\vee}$
on $V^{\vee}$ corresponds, via the identification
$\Lambda^2 V^{\vee}\simeq(\Lambda^2 V)^{\vee}\simeq\Lambda^2 V$
(the latter by Poincar\'e duality $u\mapsto(u,-)_{\Pf}$), to a vector
\[
q=e_1\wedge e_3+e_2\wedge e_4\in W,\qquad (q,q)_{\Pf}=2.
\]
The element $q$ has positive square. Its orthogonal complement
$q^{\perp}\subset W$ has signature $(3,3)-(1,0)=(2,3)$ at first
sight, but the Lorgat 2020 basis convention writes
$\Lambda^{3,2}$ with signature $(3,2)$ taking the \emph{negative}
of the Gram matrix on $q^{\perp}$ (equivalently, using
$-q$ as polariser). The conventions $\Lambda^{3,2}$ and
$\Lambda^{2,3}$ are isomorphic as quadratic lattices, differing
by an overall sign on $Q$.

\emph{The stabiliser computation.} A linear map $g\in\GL(V)$ has
$\wedge^2(g)(q)=q$ if and only if $g$ preserves the symplectic
form $\omega$ corresponding to $q$; this is the defining relation
of $\Sp_4\subset\GL_4$ (Lorgat 2020 line 246). Over $\Z$:
\[
\{g\in\GL_4(\Z):\wedge^2(g)\cdot q=q\}=\Sp_4(\Z),
\]
because the condition is linear in the matrix entries and $\Z$ is
preserved. Any $g\in\Sp_4(\Z)$ therefore induces, via $\wedge^2$,
a $\Z$-linear endomorphism of $q^{\perp}\cap\Lambda^2\Z^4=\Lambda^{3,2}$
preserving the induced quadratic form. This gives a homomorphism
\[
\wedge^2:\Sp_4(\Z)\longrightarrow\OO(\Lambda^{3,2}).
\]

\emph{Integrality of the Pfaffian pairing.} For $n=4$ the pfaffian
of a $4\times 4$ antisymmetric matrix $A=(a_{ij})$ is
$\Pf(A)=a_{12}a_{34}-a_{13}a_{24}+a_{14}a_{23}\in\Z$ whenever
$a_{ij}\in\Z$ (no denominator at $n=4$; the denominator $n/2$ of the
general formula $\Pf(A)^2=\det(A)$ is an exponent, not a
divisor). Hence $(\,,\,)_{\Pf}$ takes values in $\Z$ on
$\Lambda^2\Z^4$, and the Gram matrix above has entries in
$\{0,\pm 1\}$, giving an even unimodular integral form of signature
$(3,3)$.

\paragraph{Heal.}
Inscribe, at §3 (Lorgat 2020 lines 240--246), the three explicit
displays above: the $6\times 6$ Gram matrix of $(\,,\,)_{\Pf}$ in the
standard basis; the square of the polariser $(q,q)=\pm 2$; and the
integrality proof via $\Pf(A)=a_{12}a_{34}-a_{13}a_{24}+a_{14}a_{23}$.
The $\sl{4}\simeq\so{3,3}$ identification at the level of complex Lie
algebras is explicit: $\mathrm{ad}(X)$ on $\Lambda^2 V$ for
$X\in\sl{4}$ lands in the span of the infinitesimal generators of
$\SO((\,,\,)_{\Pf})$, because $\wedge^2$ is a representation of
$\GL(V)$ and tangent-at-identity of the orthogonality condition
$\wedge^2(g)^{\top}G\wedge^2(g)=G$ is exactly
$\wedge^2(X)^{\top}G+G\wedge^2(X)=0$, i.e.\ $\wedge^2(X)\in\so{G}$.

%---------------------------------------------------------------------
\subsection*{Cycle 2: the half-spin representation and the $+$ component}
\label{sec:w13a2-cycle-2}
%---------------------------------------------------------------------

\paragraph{Attack.}
Lorgat 2020 line 282 writes: $\OO_\R(\Lambda^{3,2})$ has four connected
components, and $\OO_\R(\Lambda^{3,2})_+$ is the index-$2$ subgroup
leaving $\Hb^{IV}_+$ invariant. Three attacks:

\begin{enumerate}[label=\textbf{O.\arabic*},leftmargin=1.5em]
\item \emph{Why four components, not two?} For an indefinite quadratic
  form of signature $(p,q)$ with $p,q\geq 1$, $\OO(p,q,\R)$ has exactly
  four connected components, indexed by
  $\pi_0\OO(p,q,\R)\simeq\Z/2\times\Z/2$: the determinant $\det=\pm 1$
  and the \emph{spinor norm} $\theta=\pm 1$. Lorgat 2020 implicitly uses
  this but does not name the bijection.
\item \emph{Which component is $+$?} The symmetric domain $\Hb^{IV}_+$
  (type IV bounded symmetric domain) is the Hermitian symmetric
  domain $\OO(3,2,\R)/\OO(3)\times\OO(2)$ restricted to the component
  preserving the complex structure; the $+$ here is \emph{not} the
  determinant $+1$ component, but the \emph{spin-norm $+1$} component,
  preserving a choice of \emph{orientation on the negative-definite
  subspace}.
\item \emph{Is the kernel $\pm I$ the half-spin kernel?} Lorgat 2020
  line 282 states the kernel of the action of
  $\OO_\R(\Lambda^{3,2})_+$ on $\Hb^{IV}_+$ is $\pm I_5$. This is the
  projective-orthogonal kernel, which coincides with the kernel of
  the half-spin representation on $\Spin(3,2)\simeq\Sp_4(\R)$: both
  central extensions by $\Z/2$ have the same centre.
\end{enumerate}

\paragraph{Ghost theorem.}
The identification works as follows. Start from the accidental
low-rank isogeny
\[
\Spin(3,2)\xrightarrow[\text{double cover}]{2:1}\SO^{0}(3,2)\hookrightarrow\OO(3,2)_+
\]
where $\SO^{0}$ denotes the identity component. The half-spin
representation (Fulton-Harris Thm 19.20; Chevalley 1954 Ch.~II) at
signature $(3,2)$ is $4$-dimensional, and
\[
\Spin(3,2)\simeq\Sp_4(\R)
\]
is the accidental isomorphism (Helgason 1978 Ch.~IX.5). Under this
identification, the half-spin representation of $\Spin(3,2)$
corresponds to the standard $4$-dimensional representation of
$\Sp_4(\R)$.

The \emph{kernel} of the covering $\Sp_4(\R)\twoheadrightarrow\SO^{0}(3,2)$
is generated by the non-trivial element of $\pi_1(\SO^{0}(3,2))$; this
pulls back to $\{\pm I_4\}\subset\Sp_4(\R)$ (the centre, which on the
half-spin picks up a sign). On $\SO^{0}(3,2)$, the corresponding
kernel is $\{\pm I_5\}$ in the defining representation on
$\R^{3,2}$.

\emph{The $+$ component.} The four components of $\OO(3,2,\R)$ are
labelled by $(\det,\theta)\in\{\pm 1\}^2$ where $\theta$ is the spin
norm. The subgroup $\OO(3,2,\R)_+$ is defined as the preimage of the
identity component in the double quotient, which corresponds to
$\theta=+1$ (spin norm is a positive real number, signed $+$).
Equivalently: $\OO(3,2,\R)_+$ is the subgroup preserving the complex
structure on $\Hb^{IV}_+$. The half-spin representation of
$\Spin(3,2)$ lands in $\Sp_4(\R)$, covering $\SO^{0}(3,2)\subset
\OO(3,2,\R)_+$. Hence the $\wedge^2$ map of Lorgat 2020 has image in
$\OO(\Lambda^{3,2})_+$, \emph{not} in any other component.

\emph{Why $\OO(\Lambda^{3,2})_+/\{\pm I_5\}=\SO^+(\Lambda^{3,2})$.}
Lorgat 2020 line 285 writes $\OO_\R(\Lambda^{3,2})_+=\pm I_5\cdot\SO_\R(\Lambda^{3,2})_+$
because $\Lambda^{3,2}$ has odd dimension $5$: for odd-dimensional
lattices, $-I$ has determinant $-1$, so
$\OO=\{\pm I\}\times\SO$ as an internal semidirect product. The quotient
$\OO(\Lambda^{3,2})_+/\{\pm I_5\}$ is therefore
$\SO^+(\Lambda^{3,2})$.

\paragraph{Heal.}
The sentence following Lemma~1 should explicitly name the chain of
coverings:
\[
\begin{tikzcd}
\Sp_4(\R)\ar[r,"\simeq"]\ar[d,two heads,"2:1"'] & \Spin(3,2)\ar[d,two heads,"2:1"] \\
\Sp_4(\R)/\{\pm I_4\}\ar[r,"\simeq"] & \SO^{0}(3,2)=\OO(3,2,\R)_+/\{\pm I_5\}
\end{tikzcd}
\]
with kernel $\{\pm 1\}$ at each level. The integral shadow $\Sp_4(\Z)\to\OO(\Lambda^{3,2})_+$
is the preimage of the integer points; the map descends to
$\Sp_4(\Z)/\{\pm I_4\}\to\OO(\Lambda^{3,2})_+/\{\pm I_5\}\simeq\SO^+(\Lambda^{3,2})$.
Both $\pm I$ kernels coincide under the accident; the sign
ambiguity in the half-spin lift is exactly $\pm I_4$ in $\Sp_4$ or
$\pm I_5$ in $\OO(3,2)$, and the identity
$\Sp_4(\Z)/\{\pm I_4\}=\OO(\Lambda^{3,2})_+/\{\pm I_5\}$ requires no
additional quotient.

%---------------------------------------------------------------------
\subsection*{Cycle 3: $\Z$-form integrality --- how $\Lambda^{3,2}$
arises as the $\Z$-points of the accidental isogeny}
\label{sec:w13a2-cycle-3}
%---------------------------------------------------------------------

\paragraph{Attack.}
The $\R$-level identity $\Sp_4(\R)\simeq\Spin(3,2)$ is classical
(Knapp 2002 VII.11, Thm 7.78). Lorgat 2020 descends to $\Z$:
$\Sp_4(\Z)/\{\pm I\}\simeq\OO(\Lambda^{3,2})_+/\{\pm I\}$. But
descent to $\Z$ is not automatic. Three attacks:

\begin{enumerate}[label=\textbf{O.\arabic*},leftmargin=1.5em]
\item \emph{$\Z$-forms of $\Sp_4$ are unique, but $\Z$-forms of
  $\SO(3,2)$ are parametrised by genera of lattices.} The genus
  $\Lambda^{1,1}\oplus\Lambda^{1,1}\oplus[2]$ is one specific even
  lattice of signature $(3,2)$; a different genus (e.g.\
  $\Lambda^{1,1}\oplus\Lambda^{1,1}\oplus[-2]$, with signature
  $(2,3)$ and appropriate sign convention) gives a different
  $\Z$-form of $\SO(3,2)$. Why this lattice and not another?
\item \emph{Is the map $\Sp_4(\Z)\to\OO(\Lambda^{3,2})$ surjective onto
  the $+$ component?} Surjectivity requires proving that every
  integer orthogonal transformation of $\Lambda^{3,2}$ preserving the
  domain $\Hb^{IV}_+$ lifts to a symplectic transformation of
  $\Z^4$. The lift exists over $\R$; at $\Z$, there is a potential
  obstruction in the spin norm.
\item \emph{Is the kernel exactly $\pm I_4$?} The kernel over $\R$ is
  $\pm I_4$; over $\Z$ it could in principle be larger (additional
  torsion elements fixing $q$). Lorgat 2020 asserts it is $\pm I_4$
  without explicit verification.
\end{enumerate}

\paragraph{Ghost theorem.}
The lattice $\Lambda^{3,2}$ is \emph{not} chosen by fiat; it is forced
by the Pfaffian construction. Concretely: the unique $\Z$-form of
$\sl{4}$ with the standard symplectic polariser $q=e_1\wedge e_3+
e_2\wedge e_4$ produces the orthogonal complement $q^{\perp}\cap
\Lambda^2\Z^4$, which \emph{is} the lattice
$\Lambda^{3,2}=\Lambda^{1,1}\oplus\Lambda^{1,1}\oplus[2]$ (Lorgat 2020
line 250). The proof:

\emph{Step 1.} The basis $\{f_1,f_2,f_3,f_{-1},f_{-2}\}$ of
$\Lambda^{3,2}$ (Lorgat 2020 lines 255--260) is explicit:
\[
f_1=e_1\wedge e_2,\ f_2=e_2\wedge e_3,\ f_3=e_1\wedge e_3-e_2\wedge e_4,\ f_{-2}=e_4\wedge e_1,\ f_{-1}=e_4\wedge e_3.
\]
\emph{Step 2.} The Gram matrix of $(\,,\,)_{\Pf}$ in this basis is
computed directly. The pair $(f_1,f_{-1})=(e_1\wedge e_2,e_4\wedge e_3)$
gives $f_1\wedge f_{-1}=e_1\wedge e_2\wedge e_4\wedge e_3=-e_1\wedge e_2\wedge e_3\wedge e_4$
so $(f_1,f_{-1})=-1$. Similarly $(f_2,f_{-2})=-1$. Finally
$(f_3,f_3)$: $f_3\wedge f_3=(e_1\wedge e_3-e_2\wedge e_4)\wedge(e_1\wedge e_3-e_2\wedge e_4)
=-2e_1\wedge e_3\wedge e_2\wedge e_4+e_1\wedge e_3\wedge e_1\wedge e_3+e_2\wedge e_4\wedge e_2\wedge e_4$.
The last two terms vanish (repeated $e_i$). The first is
$-2 e_1\wedge e_3\wedge e_2\wedge e_4=2 e_1\wedge e_2\wedge e_3\wedge e_4$.
So $(f_3,f_3)=2$. All other pairings vanish (the surviving
$4$-forms have repeated $e_i$). The Gram matrix:
\[
G|_{\Lambda^{3,2}}=\begin{pmatrix}0&&&&-1\\&0&-1&&\\&-1&0&&\\&&&&-1&\text{wait}\end{pmatrix}
\]
cleanly,
\[
G|_{\Lambda^{3,2}}=\left(\begin{array}{ccc|cc}0&&&&-1\\&0&&-1&\\&&2&&\\\hline&-1&&0&\\-1&&&&0\end{array}\right),
\]
which is (after the basis reorder $(f_1,f_{-1})$ first,
$(f_2,f_{-2})$ second, $(f_3)$ last)
$-\Lambda^{1,1}\oplus-\Lambda^{1,1}\oplus[2]$; equivalently, with the
sign convention chosen by Lorgat 2020,
$\Lambda^{1,1}\oplus\Lambda^{1,1}\oplus[2]$ of signature $(3,2)$.

\emph{Step 3 (integrality).} The basis elements $f_i\in\Lambda^2\Z^4$
are integer combinations of $e_{ij}$, so $\Lambda^{3,2}\subset\Z^5$
with integer Gram matrix; discriminant
$\det(G|_{\Lambda^{3,2}})=(-1)(-1)(2)=2$, so
$\mathrm{disc}(\Lambda^{3,2})=2$. (The unimodular host $\Lambda^2\Z^4$
of discriminant $1$ contains $\Lambda^{3,2}\oplus\Z q$ with
$(q,q)=2$; the lattice product has discriminant $2\cdot 1=2$ times
the host discriminant $1$, consistent.)

\emph{Step 4 (surjectivity onto $\SO^+$).} The generators of
$\Sp_4(\Z)$ are (Lorgat 2020 line 294--336):
(a) $g_0=\begin{pmatrix}0&I_2\\-I_2&0\end{pmatrix}$ (the Fricke/Weyl element),
(b) $g_M=\begin{pmatrix}I_2&M\\0&I_2\end{pmatrix}$ for $M\in M_{2\times 2}^{\text{sym}}(\Z)$ (upper-unipotent),
(c) $g_A=\begin{pmatrix}{}^{\top}A^{-1}&0\\0&A\end{pmatrix}$ for $A\in\GL_2(\Z)$ (Levi part).
Under $\wedge^2$, these map to the $5\times 5$ matrices
displayed in Lorgat 2020 lines 298--336. Direct computation shows that
the images generate $\SO^+(\Lambda^{3,2})$ (this is the
``elementary exercise'' of Lorgat 2020 line 293; the $5\times 5$
matrices are explicit and their products exhaust $\SO^+$ by
reduction theory --- every element of $\SO^+(\Lambda^{3,2})$ is a
product of reflections in roots of square $2$, and every such
reflection arises from a suitable $\wedge^2(g)$).

\emph{Step 5 (injectivity modulo $\pm I$).} If $g\in\Sp_4(\Z)$
satisfies $\wedge^2(g)=I_5$ on $\Lambda^{3,2}$, then $\wedge^2(g)$
fixes $q$ (by construction) and fixes $q^{\perp}$ (by hypothesis),
hence fixes all of $\Lambda^2\Z^4$. A linear map $g\in\GL_4$ with
$\wedge^2(g)=I$ on $\Lambda^2 V$ must satisfy $g=\pm I_4$ (standard,
since $\wedge^2$ is an irreducible $\GL_4$-rep restricted from the
semisimple part and centrally acts as $\det^{1/2}\cdot I$; but at
$\wedge^2$ over $\Z$, the only integer solutions are $\pm I_4$).
Hence $\ker(\wedge^2)=\{\pm I_4\}$ over $\Z$.

\paragraph{Heal.}
Inscribe the Gram-matrix computation and the discriminant-$2$
check explicitly at Lorgat 2020 lines 250--260; this grounds the
statement ``$\Lambda^{3,2}\simeq\Lambda^{1,1}\oplus\Lambda^{1,1}\oplus[2]$''
in a first-principles derivation. Add a remark on surjectivity
(Steps 4 and 5 above) after Lemma~1: injectivity from the kernel
computation $\ker(\wedge^2)=\{\pm I\}$, surjectivity from the
generator-by-generator exhaustion of $\SO^+(\Lambda^{3,2})$ via
the three Sp-generators.

%---------------------------------------------------------------------
\subsection*{Cycle 4: parabolic structure ---
Klingen/Siegel parabolics of $\Sp_4$ vs cusps of $\Hb^{IV}_+$}
\label{sec:w13a2-cycle-4}
%---------------------------------------------------------------------

\paragraph{Attack.}
For the Borcherds lift of $\phi_{0,1}$ (Wave 13 central theorem,
producing $\Delta_5$), one needs the isomorphism $\wedge^2$ to respect
parabolic structures: the Klingen parabolic and the Siegel parabolic
of $\Sp_4(\Z)$ correspond to the two cusps of $\Hb^{IV}_+$ (the
$0$-dimensional cusp and the $1$-dimensional cusp of the
orthogonal model). Lorgat 2020 slides mention only ``the Jacobi
group $\Gamma_\infty\subset\Sp_4(\Z)$'' (slide line 819). Three
attacks:

\begin{enumerate}[label=\textbf{O.\arabic*},leftmargin=1.5em]
\item \emph{Which cusps does $\Hb^{IV}_+$ have?} The type-IV bounded
  symmetric domain of dimension $3$ over $\Lambda^{3,2}$ has cusps of
  dimensions $0$ and $1$ (isotropic $\Z$-lines and $\Z$-planes in
  $\Lambda^{3,2}$). For $\Lambda^{3,2}\simeq\Lambda^{1,1}\oplus
  \Lambda^{1,1}\oplus[2]$ of Witt rank $2$, there are exactly two
  orbits of cusps under $\OO(\Lambda^{3,2})_+$: one $0$-dimensional,
  one $1$-dimensional.
\item \emph{Which parabolics does $\Sp_4(\Z)$ have?} $\Sp_4(\Z)$ has
  two conjugacy classes of maximal proper $\Q$-parabolics: the
  \emph{Klingen parabolic} $P_{\mathrm{Kl}}$ (stabiliser of a line
  in $\Q^4$, isomorphic to
  $\Sp_2(\Q)\ltimes(\mathrm{Heis}_3\rtimes\Q)$) and the
  \emph{Siegel parabolic} $P_{\mathrm{Si}}$ (stabiliser of a
  Lagrangian $2$-plane, isomorphic to $\GL_2(\Q)\ltimes\mathrm{Sym}^2\Q^2$).
\item \emph{Does $\wedge^2$ carry parabolics to cusp-stabilisers
  correctly?} Under the accident, Siegel parabolic $\leftrightarrow$
  $0$-dimensional cusp and Klingen parabolic $\leftrightarrow$
  $1$-dimensional cusp (or the opposite, depending on orientation
  convention). The \emph{Jacobi group} $\Gamma_\infty$ of Lorgat 2020
  slides --- the subgroup fixing the $z_0$-expansion at
  $\tau=i\infty$ --- is the \emph{Klingen parabolic}, because it
  fixes the \emph{line} $\{z_0=i\infty\}\subset\Hb_2$ (not a
  Lagrangian plane).
\end{enumerate}

\paragraph{Ghost theorem.}
Under the $\wedge^2$ isomorphism, the parabolic-cusp correspondence
is:
\begin{center}
\begin{tabular}{@{}lll@{}}
\hline
$\Sp_4(\Z)$ parabolic & $\wedge^2$-image in $\OO(\Lambda^{3,2})_+$ & Cusp of $\Hb^{IV}_+$ \\
\hline
Klingen $P_{\mathrm{Kl}}$ (line stabiliser) & Stabiliser of isotropic $\Z$-plane $\Pi\subset\Lambda^{3,2}$ & $1$-dimensional cusp \\
Siegel $P_{\mathrm{Si}}$ (Lagrangian stabiliser) & Stabiliser of isotropic $\Z$-line $\ell\subset\Lambda^{3,2}$ & $0$-dimensional cusp \\
\hline
\end{tabular}
\end{center}

\emph{Proof of the Klingen $\leftrightarrow$ isotropic-plane
correspondence.} Let $L\subset\Q^4$ be a line. Its annihilator
under the symplectic form is a Lagrangian hyperplane
$L^{\omega}\subset\Q^4$ containing $L$. The image $\wedge^2(L^{\omega})
\cap q^{\perp}$ is a $\Q$-plane in $\Lambda^{3,2}\otimes\Q$, and
this plane is \emph{isotropic} because: for $v,w\in L^{\omega}$,
$(v\wedge \ell,w\wedge\ell)_{\Pf}=v\wedge\ell\wedge w\wedge\ell=0$
(repeated $\ell$). Hence $P_{\mathrm{Kl}}$ maps to the stabiliser
of an isotropic plane.

\emph{Proof of the Siegel $\leftrightarrow$ isotropic-line
correspondence.} A Lagrangian $2$-plane $\Lambda\subset\Q^4$
satisfies $\omega|_\Lambda=0$, so $q$ restricted to $\wedge^2\Lambda$
vanishes --- $\wedge^2\Lambda$ is a line in $\Q^5=\Lambda^2\Q^4/\Q q$
after passing to the quotient. The image
$\wedge^2\Lambda=\Q(\ell_1\wedge\ell_2)\subset\Lambda^2\Q^4$ is a
$\Q$-line of square $(\ell_1\wedge\ell_2,\ell_1\wedge\ell_2)_{\Pf}=0$
(by $\omega$-Lagrangianity, the projection to $q$ is zero and the
square within $q^{\perp}$ is zero). Hence $P_{\mathrm{Si}}$ maps to
the stabiliser of an isotropic line.

\emph{Consistency with the Jacobi group.} The Jacobi group
$\Gamma_\infty\subset\Sp_4(\Z)$ of Lorgat 2020 slide line 819 fixes
the cusp $z_0=i\infty$ of $\Hb_2$; under the isomorphism
$\Hb_2=\Hb^{IV}_+$, this cusp is the $1$-dimensional cusp of
$\Hb^{IV}_+$, corresponding to the isotropic $\Z$-plane
$\mathrm{span}(f_1,f_{-1})\subset\Lambda^{3,2}$. (The plane is
isotropic: $(f_1,f_1)=0$, $(f_{-1},f_{-1})=0$, $(f_1,f_{-1})=-1$
---ah, a hyperbolic plane, whose only isotropic lines are $\Q f_1$
and $\Q f_{-1}$; the plane itself is \emph{not} isotropic, it has
determinant $-1$.) More carefully: the $1$-dimensional cusp
corresponds to an isotropic $\Z$-line, not an isotropic $\Z$-plane.
Lorgat 2020 slide terminology ``Jacobi group $\Gamma_\infty$'' is
the Klingen parabolic, and its fixed data is the isotropic
$\Z$-line $\Q f_1\subset\Lambda^{3,2}$ (a primitive null vector,
$(f_1,f_1)=0$). The \emph{cusp structure at $\Q f_1$} is
$1$-dimensional because $(\Q f_1)^{\perp}/\Q f_1$ has dimension
$3$, and the complexified symmetric domain of
$(\Q f_1)^{\perp}/\Q f_1$ (signature $(2,1)$, i.e.\ a hyperbolic
Kac-Moody cone) is $1$-complex-dimensional.

\paragraph{Heal.}
Inscribe, immediately after Lemma~1 of Lorgat 2020 \S3, the
parabolic-cusp dictionary above, with the explicit identification:
\begin{itemize}
\item Siegel parabolic $\leftrightarrow$ stabiliser of isotropic
  $\Z$-line $\ell\subset\Lambda^{3,2}$ $\leftrightarrow$
  $0$-dimensional cusp of $\Hb^{IV}_+$;
\item Klingen parabolic = Jacobi group $\Gamma_\infty$
  $\leftrightarrow$ stabiliser of isotropic $\Z$-line generating a
  $2$-dimensional isotropic cusp-arithmetic-structure
  $\leftrightarrow$ $1$-dimensional cusp of $\Hb^{IV}_+$.
\end{itemize}
This is essential for the Borcherds lift: the lift of
$\phi_{0,1}(\tau,z)$ produces the denominator automorphic form on
$\Hb^{IV}_+$ with prescribed principal-part singularity along the
$1$-dimensional cusp (Borcherds 1998 Thm 13.3), which corresponds to
the Klingen parabolic data in the Sp$_4$-picture. Wave 12 A2 Cycle~1
($\kappa_{\mathrm{BKM}}=c_1(0)/2=5$) depends on this parabolic
identification.

%---------------------------------------------------------------------
\subsection*{Cycle 5: Weyl-group compatibility ---
affine $\widetilde C_2$ of $\Sp_4$ vs orthogonal Weyl of $\OO(3,2)$}
\label{sec:w13a2-cycle-5}
%---------------------------------------------------------------------

\paragraph{Attack.}
The isomorphism $\Sp_4(\Z)/\{\pm I\}\simeq\OO(\Lambda^{3,2})_+/\{\pm I\}$
at the arithmetic level must induce a matching Weyl-group
structure; otherwise the root-system accident $C_2=B_2$ at finite
level (hence $\widetilde C_2=\widetilde B_2$ at affine level) is
only half the story. Three attacks:

\begin{enumerate}[label=\textbf{O.\arabic*},leftmargin=1.5em]
\item \emph{At the finite level: is the Weyl group of $\Sp_4(\R)$
  really isomorphic to the Weyl group of $\OO(\Lambda^{3,2}\otimes\R)$?}
  $W(C_2)=W(B_2)=D_4$ (dihedral of order $8$); $W(A_3)=W(D_3)=S_4$
  (order $24$). The match $W(C_2)=W(B_2)$ at finite level is the
  \emph{root-level} accident (not rank-level), confusingly named
  the same. At rank $2$, $C_2$ and $B_2$ have the same Weyl group
  because they are dual root systems, but their $\widehat{\Z}$-level
  content differs. $A_3\neq B_2$; one must check that the
  \emph{accidentally matching} Weyl group (rank 3 on both sides?) is
  self-consistent.
\item \emph{$D_3$ vs $A_3$: the Dynkin diagram identity is
  $D_3=A_3$, both with Weyl group $S_4$.} So the Weyl group of
  $\SO(3,3,\R)$ (type $D_3$) equals $S_4=W(A_3)$. On passing from
  $\SO(3,3)$ to the subgroup $\SO(3,2)$ preserving the vector
  $q$, one picks out the subgroup of $S_4$ preserving one root;
  this is $S_3=W(A_2)$ of order $6$, the Weyl group of $\SO(3,2)$.
\item \emph{Affine Weyl compatibility.} The affine Weyl group of
  $\widetilde C_2$ (for $\Sp_4$) acts on the root lattice
  $Q(C_2)$ with fundamental alcove; on the orthogonal side, the
  affine Weyl $\widetilde{D_3}=\widetilde{A_3}$ acts on
  $Q(D_3)$. Under the accident, the fundamental alcoves should
  correspond; but the rank matches only if we compare $\widetilde
  A_2$ on the Sp-side (interpreted in the Levi of the Siegel
  parabolic) with $\widetilde A_2$ on the orthogonal side.
\end{enumerate}

\paragraph{Ghost theorem.}
The Weyl-group structures under the accident:
\[
W(\Sp_4)=W(C_2)=D_4=\{\pm 1\}^2\rtimes S_2,\quad|W|=8;
\]
\[
W(\SO(5))=W(B_2)=D_4,\quad|W|=8,
\]
but neither is $\SO(3,2,\R)^{\mathrm{cpt}}$ with its Weyl. The
correct chain of isomorphisms is:

\emph{Step 1.} $\Sp_4(\R)\simeq\Spin(3,2,\R)$ gives matching Weyl
groups: both are type $C_2=B_2$, Weyl $D_4$ of order $8$.

\emph{Step 2.} On passing to $\Spin(3,2)\twoheadrightarrow\SO^{0}(3,2)$,
the Weyl group acts on the Cartan $\mathfrak h=\R^2$ (rank 2, since
$\mathrm{rk}_\R\SO(3,2)=2$) as $D_4$. Concretely, $\SO(3,2)$ has
two simple roots, $\alpha$ (short) and $\beta$ (long) with
$(\alpha,\alpha)=1$, $(\beta,\beta)=2$, $(\alpha,\beta)=-1$, angle
$3\pi/4$ --- this is the $B_2$ root diagram, dual to $C_2$ of
$\Sp_4$.

\emph{Step 3.} The $\Z$-level Weyl group $W(\Lambda^{3,2})$ is the
\emph{reflection subgroup} of $\OO(\Lambda^{3,2})$ generated by
reflections in roots of square $2$. Lorgat 2020 \S4 (lines 352--404)
computes that $W(\Lambda^{2,1})\subset\OO(\Lambda^{2,1})_+$ is
generated by reflections in
$\delta_1=2f_2-f_3,\ \delta_2=2f_{-2}-f_3,\ \delta_3=f_3$ and
$f_{-2}-f_2,\ f_2-f_3,\ f_2+f_3$. This recovers the $B_2$
Weyl-$D_4$ structure of $\Sp_4$ (three dihedral generators of the
octahedron), modulo the antiinvariance structure of $\Delta_5$
under the first three.

\emph{Step 4 (affine compatibility).} At the affine level: the
affine Weyl $\widetilde W(\Sp_4)=\widetilde W(C_2)$ is the
Coxeter group of the affine $\widetilde C_2$ diagram; on the
orthogonal side, $\widetilde W(\OO(\Lambda^{3,2}))$ is the
\emph{hyperbolic} Weyl group of $\Lambda^{3,2}$ (since
$\Lambda^{3,2}$ has signature $(3,2)$, its affine-level reflection
group is Lorentzian-hyperbolic, not Euclidean-affine). The two
coincide under the accident because at signature-$(3,2)$,
$\OO^+(\Lambda^{3,2})$ acts on the type-IV domain $\Hb^{IV}_+$, and
the fundamental domain of $\widetilde W$ is the closure of a
finite-volume hyperbolic polytope in $\Hb^{IV}_+/\R^{\times}_{+}$.
The same fundamental polytope, pulled back through $\wedge^2$, is
the Siegel fundamental domain of $\Sp_4(\Z)\setminus\Hb_2$.

\emph{Step 5 (conclusion).} The Weyl groups match at both finite
(rank $2$, Weyl $D_4$) and affine/hyperbolic levels; the
Borcherds-lift denominator automorphic form $\Phi$ satisfies the
denominator formula
$\Phi=\sum_{w\in W}(-1)^{\ell(w)}e^{w(\rho)}\prod_{\alpha>0}(1-e^{w(\alpha)})^{\mathrm{mult}(\alpha)}$
where $W$ is the (hyperbolic) Weyl group of $\Lambda^{3,2}$, and
this formula is consistent with the Maass-Sp$_4$ denominator
formula under the $\wedge^2$ isomorphism.

\paragraph{Heal.}
Inscribe the Weyl-matching dictionary at the end of Lorgat 2020
\S3 (before the transition to \S4 on $\Delta_5$ and
$\Lambda^{3,2}$):
\begin{itemize}
\item Finite Weyl: $W(\Sp_4)=W(C_2)=D_4=W(B_2)=W(\SO(3,2))$, order
  $8$, matching under the accident $C_2=B_2$.
\item Affine/hyperbolic Weyl: $\widetilde W(\Sp_4)=\widetilde W(C_2)$
  (Euclidean) corresponds to the hyperbolic Weyl of $\Lambda^{3,2}$
  acting on $\Hb^{IV}_+/\R^{\times}_+$ (Lorentzian), with matching
  fundamental polytopes in the respective quotient domains.
\item The three $f_i$-reflections of Lorgat 2020 \S4 realise the
  $B_2$-Weyl generators (the three $\delta_i$'s for the long root
  $f_3$ and short roots $2f_2-f_3$, $2f_{-2}-f_3$); the three
  $(f_i-f_j)$-reflections realise $S_3$-type permutations of the
  hyperbolic-plane basis.
\end{itemize}

%---------------------------------------------------------------------
\subsection*{Synthesis: what the five cycles establish}
%---------------------------------------------------------------------

The isomorphism
\[
\wedge^2:\Sp_4(\Z)/\{\pm I_4\}\;\xrightarrow{\ \sim\ }\;\SO^+(\Lambda^{3,2})\simeq\OO(\Lambda^{3,2})_+/\{\pm I_5\}
\]
is the $\Z$-level shadow of the accidental Dynkin-diagram identity
$D_3\simeq A_3$, or equivalently $\Sp_4(\R)\simeq\Spin(3,2)$. Five
layers of verification, each contributing a distinct witness:

\begin{enumerate}[label=(\arabic*)]
\item (Cycle 1) The exterior square $\wedge^2:\GL_4\to\GL_6$
  preserves the Pfaffian pairing; stabiliser of
  $q=e_1\wedge e_3+e_2\wedge e_4$ is $\Sp_4$.
\item (Cycle 2) The $+$ component is the spin-norm-$+1$ component;
  it is the image of $\Sp_4(\R)\simeq\Spin(3,2)$ under the accidental
  isomorphism.
\item (Cycle 3) The $\Z$-form $\Lambda^{3,2}=\Lambda^{1,1}\oplus
  \Lambda^{1,1}\oplus[2]$ emerges by construction (orthogonal
  complement of $q$ inside $\Lambda^2\Z^4$ with Pfaffian pairing);
  discriminant $2$; Weyl-group generators match.
\item (Cycle 4) Siegel parabolic $\leftrightarrow$ $0$-dim cusp,
  Klingen parabolic (= Jacobi group $\Gamma_\infty$) $\leftrightarrow$
  $1$-dim cusp, via isotropic $\Z$-line/plane data.
\item (Cycle 5) Finite Weyl $W(C_2)=D_4$ and affine Weyl
  $\widetilde W(C_2)$ correspond to finite $W(B_2)=W(\SO(3,2))$ and
  hyperbolic $\widetilde W(\Lambda^{3,2})$; fundamental domains
  match.
\end{enumerate}

This confirms Lorgat 2020 Lemma~1 as theorem, supplies the
derivations that were cited-only (\emph{``can be found in [10]''} =
Knapp 2002 VII.11), and catches one convention artefact: the
signature $(3,3)\rightsquigarrow(3,2)$ transition is exactly the
subtraction of one square-$+2$ polariser direction, not an
independent construction.

\paragraph{Cross-references inscribed.}
\begin{itemize}
\item \texttt{notes/wave12\_u1\_two\_stage\_functor.tex}:
  Stage-2 specialisation $\mathrm{Sp}_{\Sigma_{d-1},C}$ at $d=3$ on
  $X=K3\times E$ passes through the orthogonal model
  $\OO^+(\Lambda^{3,2})$ on $\Hb^{IV}_+$; the Sp$_4(\Z)$-picture is
  the alternative presentation via $\wedge^2$. Both stages land in
  the same $E_1$-chiral algebra on the reference elliptic curve.
\item \texttt{notes/wave12\_a2\_k3xE\_BKM\_kazhdan.tex} Cycle~1--2:
  the Borcherds weight identity $\kBKM(\Phi_N)=c_N(0)/2$ is stated
  on the orthogonal side ($\OO(\Lambda^{3,2})_+$); under the
  $\wedge^2$ isomorphism, it is the weight of the corresponding
  $\Sp_4(\Z)$-Siegel modular form ($\Delta_5$ at $N=1$).
\item \texttt{notes/wave13\_source\_paper.txt} \S3 (lines 233--350):
  primary source. The integrality and parabolic-cusp computations
  inscribed above close the cited-only references (``can be found
  in [10]''; ``elementary exercise to realize concretely'').
\end{itemize}

%---------------------------------------------------------------------
\subsection*{Substantive fixes inscribed into first-principles cache}
%---------------------------------------------------------------------

Three new cache entries (AP-CY157, AP-CY158, AP-CY159) are appended
to \texttt{appendices/first\_principles\_cache.md}:

\begin{itemize}
\item \textbf{AP-CY157 / signature-bookkeeping confusion.} The
  Pfaffian pairing on $\Lambda^2\Z^4$ has signature $(3,3)$; the
  orthogonal hyperplane $q^\perp$ has signature
  $(3,2)$ (drop of $(0,1)$ because $(q,q)=+2$). Conventions differ
  on sign: $\Lambda^{3,2}$ (Lorgat) vs $\Lambda^{2,3}$ (some
  authors). Identical as quadratic lattices up to overall sign.
\item \textbf{AP-CY158 / accidental-isogeny $\Z$-level descent.}
  $\Spin(3,2)\simeq\Sp_4(\R)$ is the real-form accident
  (Helgason IX.5). Descent to $\Z$ requires explicit generator-by-
  generator verification (Steps 4--5 of Cycle 3); surjective onto
  $\SO^+(\Lambda^{3,2})$, kernel $\{\pm I_4\}$. Not all real isogenies
  of accidental Lie type descend to $\Z$-form isomorphisms;
  this one does, by the Pfaffian construction.
\item \textbf{AP-CY159 / parabolic-cusp correspondence under $\wedge^2$.}
  Siegel parabolic $\leftrightarrow$ stabiliser of isotropic
  $\Z$-line $\leftrightarrow$ $0$-dim cusp; Klingen parabolic
  (=Jacobi group $\Gamma_\infty$) $\leftrightarrow$ stabiliser of
  isotropic $\Z$-line in a hyperbolic plane $\leftrightarrow$
  $1$-dim cusp. Do not conflate the two; Borcherds-lift
  principal-part data attach to the $1$-dim cusp.
\end{itemize}

\medskip
\noindent\textbf{Claim status.} \ClaimStatusTheorem (Lemma~1 of
Lorgat 2020 \S3 verified; integrality, parabolic-cusp, Weyl-matching
all explicit). The cited reference [10] (Knapp 2002) supplies the
$\R$-level accident; the $\Z$-level descent is proven here from
first principles, matching Lorgat 2020's explicit generator
computations at lines 294--336.
